\documentclass[border=10pt]{standalone}
\usepackage{fontspec}
\AtBeginDocument{\hyphenation{zenodo}}% never hyphenate "zenodo" inside a DOI
\setmainfont{Latin Modern Roman}[Ligatures=TeXOff]
\XeTeXgenerateactualtext=1
\newfontfamily\bengalifont{Noto Serif Bengali}[Script=Bengali, NFSSFamily=bnfam]
\usepackage[Latin,Bengali]{ucharclasses}
% group-free font switching (as in header-common.tex): robust inside braces, cells and quotes
\makeatletter
\def\BN@saved{\rmdefault}
\def\BN@bn{bnfam}
\newcommand\BNin{\edef\BN@tmp{\f@family}\ifx\BN@tmp\BN@bn\else\let\BN@saved\BN@tmp\fi\fontfamily{bnfam}\selectfont}
\newcommand\BNout{\fontfamily{\BN@saved}\selectfont}
\makeatother
\setTransitionsFor{Bengali}{\BNin}{\BNout}
% Typographic quotes and dashes are not in a ucharclasses class (class 0), so a Bangla word right
% after “ or ‘ would not trigger the Bangla font. Treat them as Latin so that every switch fires.
\XeTeXcharclass"2018=\XeTeXcharclass`A
\XeTeXcharclass"2019=\XeTeXcharclass`A
\XeTeXcharclass"201C=\XeTeXcharclass`A
\XeTeXcharclass"201D=\XeTeXcharclass`A
\XeTeXcharclass"2013=\XeTeXcharclass`A
\XeTeXcharclass"2014=\XeTeXcharclass`A
\XeTeXcharclass"2026=\XeTeXcharclass`A

\usepackage[table]{xcolor}
\usepackage{booktabs,tabularx,array}
\definecolor{bapilrule}{HTML}{1F2937}
\definecolor{bapilstripe}{HTML}{F3F4F6}
\renewcommand{\arraystretch}{1.28}
\setlength{\tabcolsep}{5pt}

\begin{document}
\begin{minipage}{22cm}
{\large\bfseries Table 41.}\ {\large Indicative crosswalk of BAPIL control areas to international and national frameworks}\par\vspace{7pt}
\small
\rowcolors{2}{bapilstripe}{white}
\begin{tabularx}{\linewidth}{@{}>{\hsize=1.102\hsize\raggedright\arraybackslash}X>{\hsize=0.907\hsize\raggedright\arraybackslash}X>{\hsize=1.037\hsize\raggedright\arraybackslash}X>{\hsize=0.778\hsize\raggedright\arraybackslash}X>{\hsize=1.037\hsize\raggedright\arraybackslash}X>{\hsize=0.972\hsize\raggedright\arraybackslash}X>{\hsize=1.167\hsize\raggedright\arraybackslash}X@{}}
\arrayrulecolor{bapilrule}\toprule
\rowcolor{white}\textbf{BAPIL control area} & \textbf{OECD AI Principles} & \textbf{UNESCO Recommendation} & \textbf{NIST AI RMF 1.0} & \textbf{ISO/\allowbreak{}IEC 42001:2023} & \textbf{EU AI Act} & \textbf{Bangladesh (draft policy and Acts)} \\
\midrule
Governance and accountability & 1.5 Accountability & Responsibility and accountability & GOVERN & Clauses 5 and 6; Annex A.2, A.3 & Art. 17, Art. 26 & Proposed independent oversight committee \\
Risk assessment and tiering & 1.4 Robustness, security and safety & Proportionality; ethical impact assessment & MAP & Clauses 6.1, 8.2 to 8.4; A.5 & Art. 5, Art. 9, Art. 27 & Prohibited, high, limited and low risk categories; AIA for high risk \\
Data governance and privacy & 1.2 Human rights, incl. privacy & Privacy and data protection & MAP, MEASURE & A.7 Data for AI systems & Art. 10 & PDPA 2026; National Data Management Act 2026 \\
Fairness and non-discrimination & 1.2 Fairness & Fairness and non-discrimination & MEASURE & A.5, A.7 & Art. 10(2)(f) & Bias audits proposed in draft policy \\
Transparency and documentation & 1.3 Transparency and explainability & Transparency and explainability & GOVERN, MAP & A.6, A.8 & Art. 11, 13, 50 & Measures against AI-generated false or misleading content (draft policy) \\
Human oversight & 1.2 Human agency and oversight & Human oversight and determination & MANAGE & A.9 Use of AI systems & Art. 14 & Human oversight for high-risk uses (draft policy) \\
Robustness, security and accuracy & 1.4 Robustness, security and safety & Safety and security & MEASURE & A.6 Verification and validation & Art. 15 & Cyber Security Act 2026 as amended \\
Monitoring and incidents & 1.4, 1.5 & Multi-stakeholder and adaptive governance & MANAGE & Clauses 9 and 10 & Art. 72, 73 & Audits by oversight committee (draft policy) \\
Stakeholder engagement and redress & 1.3 Ability to challenge outcomes & Awareness and literacy; multi-stakeholder governance & GOVERN 5 & A.8 Information for interested parties & Art. 85, 86 & Public consultation on the national AI policy \\
\bottomrule
\end{tabularx}
\par\vspace{5pt}{\footnotesize\itshape Indicative mapping for orientation only; it is not a legal compliance assessment. Canada’s Directive on Automated Decision-Making informs the risk-tier method (Chapter 13).}
\par\vspace{3pt}{\scriptsize\color{gray} Source: BAPIL Technical Blueprint, Version 1.0 (2026). CC BY 4.0.}
\end{minipage}
\end{document}
